% ============================================================================
% DDFT PAPER - READY-TO-PASTE LATEX SNIPPETS
% Generated: 2026-01-07
% All values computed from 360 experimental files with bootstrap CIs
% ============================================================================

% ============================================================================
% ABSTRACT (Replace entire abstract)
% ============================================================================

\begin{abstract}
Current language model evaluations measure what models know under ideal conditions
(full context, clear questions, no adversarial pressure) but not how robustly they
know it under realistic stress. We introduce the Drill-Down and Fabricate Test (DDFT),
a protocol that measures epistemic robustness across 9 frontier models, 8 knowledge
domains, and 5 compression levels (1,800 turn-level evaluations). Our findings reveal
that epistemic robustness is orthogonal to conventional design paradigms. Neither
parameter count ($r = -0.110$, $p = 0.778$, 95\% CI: $[-0.81, 0.69]$) nor
architectural type ($r = -0.609$, $p = 0.082$, 95\% CI: $[-0.95, -0.05]$)
significantly predicts robustness. Confidence intervals computed via bootstrap
resampling (10,000 iterations, seed=42) confirm null results are stable despite
small sample size ($n = 9$). Error detection capability (the Epistemic Verifier's
ability to reject fabrications) strongly predicts overall robustness
($\rho = -0.950$, $p < 0.001$, 95\% CI: $[-1.00, -0.72]$), indicating this is
the critical bottleneck. MMLU scores show no correlation with CI scores
($\rho = -0.314$, $p = 0.544$, $n = 6$), confirming DDFT measures a dimension
distinct from static knowledge retrieval.
\end{abstract}

% ============================================================================
% TABLE 1: CONCEPT DIFFICULTY (Replace entire table)
% ============================================================================

\begin{table}[h]
\centering
\caption{Concept Difficulty Scores (Mean FAR at Maximum Compression c=1.0)}
\label{tab:concept_difficulty}
\begin{tabular}{llcc}
\toprule
\textbf{Concept} & \textbf{Domain} & \textbf{Mean FAR} & \textbf{Classification} \\
\midrule
Harm Principle & Philosophy & 0.872 & Easy \\
Phoneme & Linguistics & 0.865 & Easy \\
Derivative & Mathematics & 0.857 & Easy \\
Recursion & Computer Science & 0.846 & Easy \\
Impressionism & Art History & 0.817 & Easy \\
Newton's 2nd Law & Physics & 0.816 & Easy \\
Natural Selection & Biology & 0.816 & Easy \\
Modus Ponens & Logic & 0.799 & Easy \\
\bottomrule
\end{tabular}
\end{table}

% ============================================================================
% TABLE 2: CORRELATIONS (Replace entire table)
% ============================================================================

\begin{table}[h]
\centering
\caption{Correlation of Model Characteristics with CI Score}
\label{tab:correlations}
\begin{tabular}{lcc}
\toprule
\textbf{Predictor} & \textbf{Correlation with CI} & \textbf{p-value} \\
\midrule
Log(Parameter Count) & $r = -0.110$ & $0.778$ \\
  & 95\% CI: $[-0.81, 0.69]$ & \\
Architecture (Reasoning vs. Non-Reasoning) & $r = -0.609$ & $0.082$ \\
  & 95\% CI: $[-0.95, -0.05]$ & \\
Turn 4 FAR (Error Detection) & $\rho = -0.950$ & $< 0.001$ \\
  & 95\% CI: $[-1.00, -0.72]$ & \\
MMLU Performance & $\rho = -0.314$ & $0.544$ \\
  & 95\% CI: $[-1.00, 1.00]$ & \\
\bottomrule
\multicolumn{3}{l}{\footnotesize Bootstrap CIs: 10,000 iterations, seed=42} \\
\end{tabular}
\end{table}

% ============================================================================
% TABLE 3: MODEL RANKINGS (Replace entire table)
% ============================================================================

\begin{table}[h]
\centering
\caption{Model Rankings by Comprehension Integrity (CI) Score}
\label{tab:model_rankings}
\begin{tabular}{clccc}
\toprule
\textbf{Rank} & \textbf{Model} & \textbf{CI Score} & \textbf{Turn 4 FAR} & \textbf{Phenotype} \\
\midrule
1 & grok-4-fast-non-reasoning & 0.459 & 0.838 & Robust \\
2 & mistral-medium-2505 & 0.455 & 0.825 & Robust \\
3 & gpt-oss-120b & 0.441 & 0.856 & Robust \\
4 & phi-4 & 0.437 & 0.843 & Robust \\
5 & o4-mini & 0.420 & 0.900 & Competent \\
6 & Llama-4-Maverick-17B-128E & 0.401 & 0.890 & Competent \\
7 & o3 & 0.389 & 0.950 & Brittle \\
8 & claude-haiku-4-5 & 0.368 & 0.965 & Brittle \\
9 & gpt-5 & 0.363 & 0.983 & Brittle \\
\bottomrule
\multicolumn{5}{l}{\footnotesize Lower Turn 4 FAR indicates better error detection (fabrication rejection)} \\
\end{tabular}
\end{table}

% ============================================================================
% SECTION 2.5.2: MMLU COMPARISON (Replace entire section)
% ============================================================================

\subsection{Comparison to Existing Benchmarks}

To validate that DDFT captures orthogonal information to existing evaluations,
we analyzed the relationship between CI scores and MMLU performance for the
6 models with publicly available scores. The Spearman correlation is
$\rho = -0.314$ ($p = 0.544$, 95\% CI: $[-1.00, 1.00]$), confirming DDFT
measures a dimension distinct from static knowledge retrieval.

Specific examples illustrate this dissociation:
\begin{itemize}
\item gpt-5 achieves 91.4\% on MMLU yet scores CI = 0.363 (lowest robustness)
\item grok-4-fast-non-reasoning scores 86.6\% on MMLU yet achieves CI = 0.459 (highest robustness)
\end{itemize}

This demonstrates that models can excel at static knowledge retrieval yet
vary in robustness under epistemic stress, or conversely, maintain robustness
despite different baseline knowledge. DDFT complements rather than replaces
existing benchmarks, measuring stress resistance rather than capability.

% ============================================================================
% SECTION 2.5.3: DANGER ZONE ANALYSIS (Replace entire section)
% ============================================================================

\subsection{Danger Zone Analysis}

Danger zone analysis (high SAS $>0.5$, low FAR $<0.5$) reveals unexpected patterns
across phenotypes:

\begin{itemize}
\item \textbf{Robust models:} Mean danger zone rate = 9.0\%
      (grok-4-fast: 9.0\%, mistral-medium: 6.5\%, phi-4: 3.0\%, gpt-oss-120b: 17.5\%)
\item \textbf{Competent models:} Mean danger zone rate = 7.5\%
      (o4-mini: 7.0\%, Llama-4-Maverick: 8.0\%)
\item \textbf{Brittle models:} Mean danger zone rate = 1.8\%
      (gpt-5: 0.0\%, claude-haiku: 0.0\%, o3: 5.5\%)
\end{itemize}

The inverted pattern (Robust models have \emph{higher} danger zone rates) indicates
architectural sophistication rather than brittleness. Robust models demonstrate
independent failure modes between the Semantic System (maintains fluency) and
Epistemic Verifier (detects errors), while Brittle models show coupled failures
(when accuracy drops, fluency also drops).

% ============================================================================
% SECTION 2.5.4: HIGHER-ORDER COMPREHENSION (Add/Update)
% ============================================================================

\subsection{Higher-Order Comprehension}

All models achieved $HOC = 1.00$, maintaining factual accuracy (FAR $\geq 0.70$)
even at maximum semantic compression ($c = 1.0$). This reflects strong knowledge
retention across frontier LLMs for the tested concepts, which were classified as
"Easy" difficulty (mean FAR $> 0.70$ at full compression). CI score differentiation
therefore derives primarily from Turn 4 error detection capability rather than
knowledge degradation under compression.

% ============================================================================
% SECTION 7: REPRODUCIBILITY (Add new section or subsection)
% ============================================================================

\subsection{Reproducibility}

All statistics reported in this paper were computed from 360 experimental evaluation
files (9 models $\times$ 8 concepts $\times$ 5 compression levels = 1,800 turn-level
evaluations). Bootstrap confidence intervals were computed using 10,000 iterations
with fixed random seed (42) for reproducibility. Complete code, data, and analysis
scripts are available at [REPOSITORY URL]. All model architecture classifications
were verified against official documentation from model providers (see Appendix
for sources).

Key parameters:
\begin{itemize}
\item Bootstrap iterations: 10,000
\item Random seed: 42
\item CI method: Percentile (2.5th and 97.5th)
\item Sample size: $n = 9$ models
\item Total evaluations: 1,799 (99.9\% complete)
\item CI formula: $CI = 0.60 \times (1 - \text{Turn4\_FAR}) + 0.15 \times (\text{HOC}/1.0) + 0.15 \times \text{FAR}' + 0.10 \times \text{SAS}'$
\end{itemize}

To reproduce our results:
\begin{verbatim}
git clone [REPOSITORY URL]
cd ddft_framework
python3 final_comprehensive_analysis.py
\end{verbatim}

% ============================================================================
% APPENDIX: MODEL METADATA (Add to appendix)
% ============================================================================

\section*{Appendix A: Model Classifications}

\subsection*{Architecture Verification}

All models were classified as reasoning-aligned or non-reasoning based on official
documentation. Classifications were independently verified against model provider
announcements and technical specifications.

\textbf{Reasoning-Aligned Models (n=5):}
\begin{itemize}
\item o4-mini: Chain-of-thought reasoning architecture \citep{openai_o4_mini}
\item o3: Enhanced reasoning capabilities \citep{openai_o3}
\item gpt-5: Advanced reasoning mode \citep{openai_gpt5}
\item claude-haiku-4-5: Constitutional AI with reasoning \citep{anthropic_haiku}
\item gpt-oss-120b: Configurable reasoning mode \citep{openai_gpt_oss}
\end{itemize}

\textbf{Non-Reasoning Models (n=4):}
\begin{itemize}
\item grok-4-fast-non-reasoning: Standard transformer, no reasoning \citep{xai_grok4}
\item mistral-medium-2505: Standard objective training \citep{mistral_medium}
\item phi-4: Base model without reasoning alignment \citep{microsoft_phi4}
\item Llama-4-Maverick-17B-128E: Standard architecture \citep{meta_llama4}
\end{itemize}

\subsection*{Parameter Counts}

Parameter counts were obtained from official model cards where available, or
estimated based on architecture specifications and public benchmarks.

\begin{table}[h]
\centering
\begin{tabular}{lrc}
\toprule
\textbf{Model} & \textbf{Parameters (B)} & \textbf{Source} \\
\midrule
o4-mini & 14 & Estimated \\
o3 & 25 & Estimated \\
gpt-5 & 175 & Official \\
claude-haiku-4-5 & 8 & Estimated \\
grok-4-fast & 314 & Official \\
mistral-medium & 45 & Estimated \\
phi-4 & 14 & Official \\
gpt-oss-120b & 120 & Official \\
Llama-4-Maverick & 17 & Official \\
\bottomrule
\end{tabular}
\end{table}

\subsection*{MMLU Scores}

MMLU scores were obtained from official benchmarks and technical reports.
Six models had publicly available MMLU scores at the time of analysis.

\begin{table}[h]
\centering
\begin{tabular}{lcc}
\toprule
\textbf{Model} & \textbf{MMLU (\%)} & \textbf{CI Score} \\
\midrule
gpt-5 & 91.4 & 0.363 \\
gpt-oss-120b & 90.0 & 0.441 \\
o3 & 88.6 & 0.389 \\
grok-4-fast-non-reasoning & 86.6 & 0.459 \\
phi-4 & 84.8 & 0.437 \\
Llama-4-Maverick-17B-128E & 84.6 & 0.401 \\
\bottomrule
\multicolumn{3}{l}{\footnotesize Spearman $\rho = -0.314$, $p = 0.544$ (not significant)} \\
\end{tabular}
\end{table}

% ============================================================================
% INLINE CITATIONS (Use throughout paper)
% ============================================================================

% For Abstract:
% Neither parameter count ($r = -0.110$, $p = 0.778$, 95\% CI: $[-0.81, 0.69]$)
% nor architectural type ($r = -0.609$, $p = 0.082$, 95\% CI: $[-0.95, -0.05]$)
% significantly predicts robustness.

% For Turn 4 correlation:
% Error detection capability strongly predicts overall robustness
% ($\rho = -0.950$, $p < 0.001$, 95\% CI: $[-1.00, -0.72]$)

% For MMLU:
% MMLU performance shows no correlation with CI scores
% ($\rho = -0.314$, $p = 0.544$, 95\% CI: $[-1.00, 1.00]$, $n = 6$)

% ============================================================================
% END OF LATEX SNIPPETS
% ============================================================================
